% GENERATED FILE. Edit canonical lessons, then run npm run generate:books.
% canonical-source-hash: fnv1a64:51e471765d8e84fe
% canonical-lessons: FA-C464-maghaze, FA-C464-maghazedar, FA-C464-moamele-kardan, FA-C464-beyane, FA-C464-geran-forukhtan

\chapter{Shops and deals}
\label{ch:fa-shops-and-deals}
\hlchaptermodality{464}

\begin{chapteropening}
\textbf{By the end of this chapter:} \emph{I can deal with shopkeepers: a shop, a shopkeeper, doing a deal, a deposit and being overcharged.}

Complete the last lesson of chapter 464: I can deal with shopkeepers: a shop, a shopkeeper, doing a deal, a deposit and being overcharged.
\end{chapteropening}

\section[\texorpdfstring{\fa{مغازه} (maghâze)}{مغازه (maghâze)}]{\fa{مغازه} (maghâze) — a shop (US: a store)}

\label{lesson:FA-C464-maghaze}



\begin{quote}
Before the new one: say the Persian for jewels, then the Persian for a perfume.
\end{quote}

\subsection*{You'll want to know: \fa{مغازه}}
\textbf{\fa{مغازه}} (\emph{maghâze}) — \textquotedblleft{}a shop (US: a store)\textquotedblright{}.

\begin{quote}
\fa{این} \fa{مغازه} \fa{تا} \fa{ساعت} \fa{ده} \fa{شب} \fa{باز} \fa{است}.
\end{quote}

\begin{quote}
\emph{In maghâze tâ sâ'at-e dah-e shab bâz ast.}
\end{quote}

\begin{quote}
This shop is open until ten at night.
\end{quote}

The everyday word; \emph{forushgâh} is a bigger store.

\begin{cousinweb}
\emph{maghâze} was borrowed from Russian \emph{magazín} 'shop' (EtymDB), which comes in the end from Arabic \emph{maḫāzin} 'storehouses'. French \emph{magasin} and English \emph{magazine} go back to the same Arabic word.
\end{cousinweb}

\begin{grammarlens}[title={What you've built}]
One more word for telling, explaining and describing at B1.
\end{grammarlens}

\subsection*{Your turn}
\begin{itemize}
\raggedright
  \item \emph{Say it:} \emph{maghâze}
  \item \emph{Say it:} \emph{maghâze}, once more
  \item \emph{Say it:} a sentence of your own with \emph{maghâze}
  \item \emph{Recall:} say the Persian for jewels, then the Persian for a perfume, then say \emph{maghâze} again
\end{itemize}

\begin{tcolorbox}[breakable,colback=teal!4,colframe=teal!35!black,title={Before you move on}]
What does \fa{مغازه} mean? (\textquotedblleft{}A shop (US: a store)\textquotedblright{}.) Say it once more.
\end{tcolorbox}
\section[\texorpdfstring{\fa{مغازه} \fa{دار} (maghâzedâr)}{مغازه دار (maghâzedâr)}]{\fa{مغازه} \fa{دار} (maghâzedâr) — a shopkeeper}

\label{lesson:FA-C464-maghazedar}



\begin{quote}
Before the new one: say the Persian for a perfume, then the Persian for a shop.
\end{quote}

\subsection*{You'll want to know: \fa{مغازه} \fa{دار}}
\textbf{\fa{مغازه} \fa{دار}} (\emph{maghâzedâr}) — \textquotedblleft{}a shopkeeper\textquotedblright{}.

\begin{quote}
\fa{مغازه} \fa{دار} \fa{به} \fa{ما} \fa{تخفیف} \fa{داد}.
\end{quote}

\begin{quote}
\emph{Maghâzedâr be mâ takhfif dâd.}
\end{quote}

\begin{quote}
The shopkeeper gave us a discount.
\end{quote}

\emph{maghâze} 'shop' + \emph{dâr} 'keeper'.

\begin{cousinweb}
\emph{maghâzedâr} is built from \emph{maghâze} 'shop', which EtymDB takes from Russian \emph{magazín}, and native \emph{dâr} 'holder, keeper', from the stem of \emph{dâshtan} 'to have'.
\end{cousinweb}

\begin{grammarlens}[title={What you've built}]
One more word for telling, explaining and describing at B1.
\end{grammarlens}

\subsection*{Your turn}
\begin{itemize}
\raggedright
  \item \emph{Say it:} \emph{maghâzedâr}
  \item \emph{Say it:} \emph{maghâzedâr}, once more
  \item \emph{Say it:} a sentence of your own with \emph{maghâzedâr}
  \item \emph{Recall:} say the Persian for a perfume, then the Persian for a shop, then say \emph{maghâzedâr} again
\end{itemize}

\begin{tcolorbox}[breakable,colback=teal!4,colframe=teal!35!black,title={Before you move on}]
What does \fa{مغازه} \fa{دار} mean? (\textquotedblleft{}A shopkeeper\textquotedblright{}.) Say it once more.
\end{tcolorbox}
\section[\texorpdfstring{\fa{معامله} \fa{کردن} (mo'âmele kardan)}{معامله کردن (mo'âmele kardan)}]{\fa{معامله} \fa{کردن} (mo'âmele kardan) — to do a deal, to trade}

\label{lesson:FA-C464-moamele-kardan}



\begin{quote}
Before the new one: say the Persian for a shop, then the Persian for a shopkeeper.
\end{quote}

\subsection*{You'll want to know: \fa{معامله} \fa{کردن}}
\textbf{\fa{معامله} \fa{کردن}} (\emph{mo'âmele kardan}) — \textquotedblleft{}to do a deal, to trade\textquotedblright{}.

\begin{quote}
\fa{پدرم} \fa{سالها} \fa{با} \fa{این} \fa{فروشنده} \fa{معامله} \fa{کرده} \fa{است}.
\end{quote}

\begin{quote}
\emph{Pedaram sâl-hâ bâ in forushande mo'âmele karde ast.}
\end{quote}

\begin{quote}
My father has done business with this seller for years.
\end{quote}

\emph{mo'âmele} 'a deal' + \emph{kardan}; the other side takes \emph{bâ}.

\begin{cousinweb}
\emph{mo'âmele kardan} is built from \emph{mo'âmele} 'a deal', borrowed from Arabic \emph{muʿāmala} 'conduct, dealing' (EtymDB), and \emph{kardan} 'to do', inherited from Middle Persian \emph{kardan} (EtymDB).
\end{cousinweb}

\begin{grammarlens}[title={What you've built}]
One more word for telling, explaining and describing at B1.
\end{grammarlens}

\subsection*{Your turn}
\begin{itemize}
\raggedright
  \item \emph{Say it:} \emph{mo'âmele kardan}
  \item \emph{Say it:} \emph{mo'âmele kardan}, once more
  \item \emph{Say it:} a sentence of your own with \emph{mo'âmele kardan}
  \item \emph{Recall:} say the Persian for a shop, then the Persian for a shopkeeper, then say \emph{mo'âmele kardan} again
\end{itemize}

\begin{tcolorbox}[breakable,colback=teal!4,colframe=teal!35!black,title={Before you move on}]
What does \fa{معامله} \fa{کردن} mean? (\textquotedblleft{}To do a deal, to trade\textquotedblright{}.) Say it once more.
\end{tcolorbox}
\section[\texorpdfstring{\fa{بیعانه} (bey'âne)}{بیعانه (bey'âne)}]{\fa{بیعانه} (bey'âne) — a down payment (to hold a purchase); a deposit}

\label{lesson:FA-C464-beyane}



\begin{quote}
Before the new one: say the Persian for a shopkeeper, then the Persian for to do a deal.
\end{quote}

\subsection*{You'll want to know: \fa{بیعانه}}
\textbf{\fa{بیعانه}} (\emph{bey'âne}) — \textquotedblleft{}a down payment (to hold a purchase); a deposit\textquotedblright{}.

\begin{quote}
\fa{برای} \fa{فرش} \fa{بیعانه} \fa{دادیم} \fa{و} \fa{فردا} \fa{بقیه} \fa{پول} \fa{را} \fa{میدهیم}.
\end{quote}

\begin{quote}
\emph{Barâ-ye farsh bey'âne dâdim va fardâ baqiye-ye pul râ midahim.}
\end{quote}

\begin{quote}
We paid a deposit on the carpet and we'll pay the rest tomorrow.
\end{quote}

Paid when you agree to buy: a house, a car, a carpet.

\begin{cousinweb}
\emph{bey'âne} is usually traced to Arabic \emph{bayʿ} 'sale', with the Persian ending \emph{-âne}. EtymDB does not list it.
\end{cousinweb}

\begin{grammarlens}[title={What you've built}]
One more word for telling, explaining and describing at B1.
\end{grammarlens}

\subsection*{Your turn}
\begin{itemize}
\raggedright
  \item \emph{Say it:} \emph{bey'âne}
  \item \emph{Say it:} \emph{bey'âne}, once more
  \item \emph{Say it:} a sentence of your own with \emph{bey'âne}
  \item \emph{Recall:} say the Persian for a shopkeeper, then the Persian for to do a deal, then say \emph{bey'âne} again
\end{itemize}

\begin{tcolorbox}[breakable,colback=teal!4,colframe=teal!35!black,title={Before you move on}]
What does \fa{بیعانه} mean? (\textquotedblleft{}A down payment (to hold a purchase); a deposit\textquotedblright{}.) Say it once more.
\end{tcolorbox}
\section[\texorpdfstring{\fa{گران} \fa{فروختن} (gerân forukhtan)}{گران فروختن (gerân forukhtan)}]{\fa{گران} \fa{فروختن} (gerân forukhtan) — to overcharge, to sell at too high a price}

\label{lesson:FA-C464-geran-forukhtan}



\begin{quote}
Before the new one: say the Persian for to do a deal, then the Persian for a down payment.
\end{quote}

\subsection*{You'll want to know: \fa{گران} \fa{فروختن}}
\textbf{\fa{گران} \fa{فروختن}} (\emph{gerân forukhtan}) — \textquotedblleft{}to overcharge, to sell at too high a price\textquotedblright{}.

\begin{quote}
\fa{این} \fa{فروشنده} \fa{به} \fa{مسافرها} \fa{گران} \fa{میفروشد}.
\end{quote}

\begin{quote}
\emph{In forushande be mosâfer-hâ gerân miforushad.}
\end{quote}

\begin{quote}
This seller overcharges travellers.
\end{quote}

\emph{gerân} 'dear' + \emph{forukhtan} 'to sell'. Present stem \emph{forush-}: \emph{gerân miforushad}.

\begin{cousinweb}
\emph{gerân forukhtan} is built from \emph{gerân} 'dear, expensive', which EtymDB traces through Middle Persian \emph{grān} to PIE \emph{g\textsuperscript{w}réh\textsubscript{2}us} 'heavy', and \emph{forukhtan} 'to sell', usually traced to Middle Persian \emph{frōxtan}.
\end{cousinweb}

\begin{grammarlens}[title={What you've built}]
One more word for telling, explaining and describing at B1.
\end{grammarlens}

\subsection*{Your turn}
\begin{itemize}
\raggedright
  \item \emph{Say it:} \emph{gerân forukhtan}
  \item \emph{Say it:} \emph{gerân forukhtan}, once more
  \item \emph{Say it:} a sentence of your own with \emph{gerân forukhtan}
  \item \emph{Recall:} say the Persian for to do a deal, then the Persian for a down payment, then say \emph{gerân forukhtan} again
\end{itemize}

\begin{tcolorbox}[breakable,colback=teal!4,colframe=teal!35!black,title={Before you move on}]
What does \fa{گران} \fa{فروختن} mean? (\textquotedblleft{}To overcharge, to sell at too high a price\textquotedblright{}.) Say it once more.
\end{tcolorbox}
